% Hafta 1 — Belleğin dört ömrü ve her adımın CWE'si
% Derleme: py -3.12 tools/latex/derle.py docs/week-1/assets/h01-09-bellek-omru.tex
\documentclass[tikz,border=8pt]{standalone}
\usepackage{cen429-cizim}
\begin{document}
\begin{tikzpicture}[node distance=24mm and 24mm]
  \node[baslik] (b) at (0,0) {Belleğin dört ömrü ve her adımın CWE'si};
  \node[altbaslik] (alt) at (b.south west) {her adımın atlanması ayrı bir zafiyet sınıfı üretir};

  \node[kutu=32mm, below=8mm of alt.south west, anchor=north west, xshift=2mm] (k1) {\kb{1 · AYIR}\\malloc / new};
  \node[uyarikutu=32mm, right=of k1] (k2) {\kb{2 · DENETLE}\\NULL mı?\\CWE-476};
  \node[kutu=32mm, right=of k2] (k3) {\kb{3 · KULLAN}\\sınır içinde\\CWE-787};

  \draw[ok, draw=soluk] (k1.east) -- (k2.west);
  \draw[ok, draw=soluk] (k2.east) -- (k3.west);

  \node[iyikutu=32mm, below=24mm of k1] (k4) {\kb{4 · SİL}\\sırrı temizle\\CWE-226};
  \node[iyikutu=32mm, right=of k4] (k5) {\kb{5 · SERBEST BIRAK}\\bir kez\\CWE-415};
  \node[morkutu=32mm, right=of k5] (k6) {\kb{6 · İŞARETÇİYİ UNUT}\\p = NULL\\CWE-416};

  \draw[ok, draw=soluk] (k4.east) -- (k5.west);
  \draw[ok, draw=soluk] (k5.east) -- (k6.west);
  \draw[ok, draw=soluk] (k3.south) -- ++(0,-8mm) -| (k4.north);

  \node[dip, text width=140mm, below=9mm of k6.south -| k3, anchor=north]
    {Altı adımın hepsi yapılmazsa bellek hatası kaçınılmazdır.};
\end{tikzpicture}
\end{document}
